\documentclass{report}
\usepackage{amsmath,amssymb}
\setlength{\parindent}{0mm}
\setlength{\parskip}{1em}
\begin{document}
\begin{center}
\rule{6in}{1pt} \
{\large Karen Willcox \\
{\bf Model Reduction for Uncertainty Quantification of Large-scale Systems}}

77 Massachusetts Avenue \\ Room 33-207 \\ Cambridge \\ MA 02139 \\ USA
\\
{\tt kwillcox@mit.edu}\\
Phuong Hunyh\\
Chad Lieberman\\
Leo Ng\end{center}

Uncertainty quantification is becoming recognized as an essential aspect
of development and use of numerical simulation tools, yet it remains
computationally intractable for large-scale complex systems characterized
by high-dimensional uncertainty spaces. For example, efficient numerical
tools to support decision under uncertainty are essential in many
settings. Examples include the design, optimization and control of
complex engineered systems, and the real-time characterization of and
response to hazardous events. In both these settings, it is essential to
address the many different sources and types of uncertainty, including
those stemming from the fidelity of the underlying mathematical models.
However, solving the resulting stochastic optimization problem may be
computationally infeasible, especially when the underlying simulation
models result from discretization of systems of partial differential
equations. Model reduction thus has an important role to play in
producing low-order approximate models that retain the essential system
dynamics but that are fast to solve.

Reduced-order models are commonly derived using a projection framework in
which the governing equations of the forward model are projected onto a
subspace of reduced dimension. This reduced
subspace is defined via a set of basis vectors. Methods to compute the
basis in the large-scale setting include approximate balanced truncation,
Krylov-subspace methods, proper orthogonal decomposition (POD), and
reduced basis methods. Balanced truncation and Krylov-subspace methods
are in general restricted to linear systems. POD and reduced basis
methods use empirical information generated from sampled solutions
(``snapshots") of the large-scale system to create the reduced basis.

This talk will discuss formulation and derivation of projection-based
reduced-order models for uncertainty quantification of systems governed
by partial differential equations. In particular, we demonstrate the use
of reduced models for uncertainty propagation, solution of statistical
inverse problems, and optimization under uncertainty. Our methods use
state approximations through the POD, reductions in parameter
dimensionality through parameter basis approximations, and the discrete
empirical interpolation method for efficient evaluation of nonlinear
terms.


\end{document}
